\section*{Chapitre \ref{ch:in:reading-a-trading-firms-accounts} --- Lecture des comptes d'une société de trading}

\begin{solution}{exo:in:reading-a-trading-firms-accounts:1}
Coût du personnel par salarié~: $556.1/173=\pounds3.21$ millions de livres~; chiffre d'affaires par salarié~: $1\,224.8/173=\pounds7.08$ millions de livres~; \omterm{def:in:reading-a-trading-firms-accounts:staff}{charges de personnel}
égales à 45.4~\% du chiffre d'affaires.
\end{solution}

\begin{solution}{exo:in:reading-a-trading-firms-accounts:2}
$434.9/483=0.90$ millions de dollars pour l'exercice de 14 mois~; sur douze mois, $0.90\times12/14=0.77$ millions de dollars. La première période
d'une nouvelle entité démarre aussi avec un effectif partiel, de sorte que même le chiffre annualisé n'est qu'indicatif.
\end{solution}

\begin{solution}{exo:in:reading-a-trading-firms-accounts:3}
Les tranches totalisent 717 personnes~; $312+150=462$ d'entre elles, soit 64.4~\%, ont reçu moins de 2~millions d'euros.
\end{solution}

\begin{solution}{exo:in:reading-a-trading-firms-accounts:4}
Avec le signe porté par la balise, $1\,224.817-686.012-17.751=\pounds521.054$ millions de livres, contre 556.556~millions de livres annoncés~: un écart
de 35.502~millions de livres, soit exactement le double des autres produits. Avec $+17.751$, les lignes s'additionnent, donc c'est le signe de la balise qui est faux,
et non l'état lui-même. Un programme doit conserver le fait tel qu'il est balisé, signaler le rapprochement échoué avec ses deux valeurs, et
n'appliquer une correction que comme une surcharge documentée qui conserve l'original.
\end{solution}

\begin{solution}{exo:in:reading-a-trading-firms-accounts:5}
Au taux de 1.2785~\$ pour une livre, $3.21\times1.2785=4.11$ millions de dollars, contre 1.25~millions de dollars (2024) et 2.11~millions de dollars (2025) dans
la société de personnes. Comparaison peu probante, car~: les sociétés diffèrent par leur activité et leur taille~; les parts de bénéfice des associés, supérieures
à ses \omterm{def:in:reading-a-trading-firms-accounts:staff}{charges de personnel}, ne sont pas dans la part retenue~; les exercices comptables ne coïncident pas~; et les deux sont des moyennes sur des rémunérations asymétriques.
\end{solution}

\begin{solution}{exo:in:reading-a-trading-firms-accounts:6}
$1\,855.4/8=232$ millions de dollars par associé ne décrit personne~: les associés comprennent des personnes morales, et celui qui a le
droit d'accès le plus élevé, une personne morale, a reçu 84.3~\% du bénéfice. Cette part sort du groupe~; ce que reçoivent les
associés pris individuellement n'apparaît pas dans les comptes.
\end{solution}

\begin{solution}{exo:in:reading-a-trading-firms-accounts:7}
Ratio de rémunération rapporté au chiffre d'affaires~: 12.1~\%, 13.4~\%, 16.5~\%, 17.2~\%, 15.1~\%, 14.5~\% (2020 à 2025)~; rapporté aux revenus de trading~: 15.8~\%,
17.9~\%, 24.0~\%, 30.3~\%, 23.9~\%, 21.7~\%. Le ratio rapporté aux revenus de trading varie davantage, car ces revenus ont presque diminué de
moitié de 2020 à 2023, alors que la rémunération est restée proche de 0.39~milliard de dollars~; le chiffre d'affaires, qui inclut d'autres
lignes que les revenus de trading, a moins baissé.
\end{solution}

\begin{solution}{exo:in:reading-a-trading-firms-accounts:8}
Le chiffre d'affaires de la société est la part des revenus de trading du groupe que le prix de transfert lui assigne, et ses cinq salariés
ne sont pas les traders~: depuis 2022, le personnel londonien, 790 en moyenne en 2025, est employé par la société de personnes et détaché
selon les besoins. Le chiffre d'affaires par salarié d'une entité qui n'emploie pas les personnes qui génèrent ce chiffre d'affaires ne mesure rien.
Les \omterm{def:in:reading-a-trading-firms-accounts:consolidated}{comptes consolidés} du groupe, qui, eux, le feraient, ne figurent pas parmi les dépôts britanniques.
\end{solution}

\begin{solution}{pb:in:reading-a-trading-firms-accounts:1}
\begin{enumerate}
\item Un registre public auprès duquel les sociétés et les sociétés de personnes déposent leurs statuts, dirigeants, propriétaires et comptes~;
les comptes annuels que la loi oblige une entité à établir, à faire certifier sauf dispense, et à déposer.
\item \omterm{def:in:reading-a-trading-firms-accounts:registry}{Registres des sociétés}~; SEC EDGAR~; rapports annuels des sociétés cotées~; prospectus et documents d'offre~; rapports des banques
par branche d'activité et \omterm{def:in:reading-a-trading-firms-accounts:pillar3}{publications au titre du Pilier~3}.
\item Environ 75~\% des quelque 2.2~millions de dépôts par an arrivent sous forme électronique (\omterm{def:in:reading-a-trading-firms-accounts:ixbrl}{Inline XBRL} ou XBRL)~; le reste se compose de documents papier ou
images, sans balises.
\item Le rapport périodique Form X-17A-5 sur la situation financière, le résultat et les fonds propres nets d'un broker-dealer, confidentiel~;
l'état annuel certifié de la situation financière, déposé comme document distinct, est public.
\item Le rapport public d'une banque ou d'une grande entreprise d'investissement sur les fonds propres, les risques et la rémunération~; il recense les personnes
rémunérées à 1~million d'euros ou davantage, par tranches.
\item Les charges de l'ensemble des personnes employées au cours de l'année~; la moyenne mensuelle des personnes employées sous
contrat de travail~; les salaires et traitements, les charges sociales, les autres coûts de retraite.
\item Les prestataires, les salariés détachés par d'autres sociétés du groupe, et les associés ou membres.
\item Une société mère et ses filiales présentées comme une seule entité, les opérations intragroupe étant éliminées~; ce que touchent les associés
d'une société de personnes à responsabilité limitée (LLP)~: montants salariés et parts de bénéfice automatiques, comptabilisés en charges, plus le bénéfice réparti à leur
gré.
\item XHTML lisible par l'humain dont les chiffres sont balisés par des concepts de taxonomie~; le contexte donne l'entité, la période et
les dimensions, l'unité (la devise ou la mesure), l'échelle (la puissance de dix), le signe, le format des séparateurs.
\item Chaque période réapparaît dans tous les dépôts ultérieurs comme comparatif, et un dépôt ultérieur peut la retraiter.
\item Chiffre d'affaires~: 140.6, 269.8, 778.6, 664.2, 588.2, 1\,224.8 millions de livres~; \omterm{def:in:reading-a-trading-firms-accounts:staff}{charges de personnel}~: 116.7, 182.9, 409.0, 394.7,
408.0, 556.1 millions de livres~; effectifs~: 68, 83, 98, 113, 143, 173.
\item \omterm{def:in:reading-a-trading-firms-accounts:staff}{Charges de personnel} par salarié~: 1.72, 2.20, 4.17, 3.49, 2.85, 3.21 millions de livres~; chiffre d'affaires par salarié~: 2.07, 3.25, 7.94,
5.88, 4.11, 7.08 millions de livres.
\item 83.0~\%, 67.8~\%, 52.5~\%, 59.4~\%, 69.4~\%, 45.4~\%~: le plus élevé pour l'exercice comptable 2020, lorsque le chiffre d'affaires était le plus faible, et le plus bas pour
l'exercice comptable 2025, lorsqu'il était le plus élevé~; une part de la rémunération ne suit pas le chiffre d'affaires.
\item Chiffre d'affaires moins frais administratifs plus autres produits, avec le signe porté par la balise, donne 521.1~millions de livres contre le
montant annoncé de 556.6~millions de livres~; les autres produits sont balisés en négatif alors que ce sont bien des produits.
\item Gains de fair value sur des placements inscrits sous le résultat d'exploitation~: 233.9, 249.7, 350.7 et 178.3 millions de livres.
\item Dividendes de 1\,694.4~millions de livres contre un bénéfice après impôt de 1\,718.4~millions de livres~: 98.6~\%.
\item \omterm{def:in:reading-a-trading-firms-accounts:staff}{Charges de personnel} par salarié de 1.72 à 4.17~millions de livres et chiffre d'affaires par salarié de 2.07 à 7.94~millions de livres sur les exercices comptables 2020 à 2025~;
98.6~\% du bénéfice après impôt de six ans ont été distribués.
\item C'est une moyenne sur une distribution asymétrique~; elle inclut les charges sociales et de retraite de l'employeur~; elle mélange
les rémunérations différées des années précédentes avec le coût de cette année.
\item Une société de personnes sort le bénéfice de ses associés des \omterm{def:in:reading-a-trading-firms-accounts:staff}{charges de personnel} et peut compter des personnes morales parmi ses associés~; le chiffre d'affaires d'une entité du groupe
est fixé par le prix de transfert, et son personnel peut être employé ailleurs dans le groupe.
\item Ils lui indiquent la taille de la société, la manière dont ses coûts et ses rémunérations évoluent avec son chiffre d'affaires, et qu'elle distribue presque tout son
bénéfice~; ils ne lui disent pas ce qu'elle, ou un salarié type de son poste, percevra.
\end{enumerate}
\end{solution}

\begin{solution}{iq:in:reading-a-trading-firms-accounts:1}
La valeur telle que balisée~: 900, sans signe moins si la balise n'en a pas~; les parenthèses relèvent de la présentation. L'attribut de signe de
la balise et le solde du concept (débit ou crédit) dans la taxonomie donnent le sens.
\iqlookfor{distinguer présentation et données~; savoir que les charges sont généralement balisées en positif.}
\end{solution}

\begin{solution}{iq:in:reading-a-trading-firms-accounts:2}
Les \omterm{def:in:reading-a-trading-firms-accounts:consolidated}{comptes consolidés} additionnent ligne à ligne la société mère et ses filiales et éliminent les opérations intragroupe~; les comptes de la société mère
présentent la seule société mère, les filiales figurant comme participations. Comparer des \omterm{def:in:reading-a-trading-firms-accounts:consolidated}{comptes consolidés} à des \omterm{def:in:reading-a-trading-firms-accounts:consolidated}{comptes consolidés}, à la même
date et dans une seule devise.
\iqlookfor{élimination des opérations intragroupe~; comparer ce qui est comparable.}
\end{solution}

\begin{solution}{iq:in:reading-a-trading-firms-accounts:3}
Télécharger chaque fichier quotidien, lire ses documents en flux, conserver ceux des sociétés figurant dans une liste de sociétés de trading (par code SIC et une
liste établie à la main), analyser les faits, et stocker chacun avec la société, la période, le concept, la valeur, le fichier source et la date de dépôt.
Conserver chaque version~: un dépôt ultérieur qui retraite une période forme une nouvelle ligne, et les requêtes retiennent la plus récente. Rapprocher les totaux
et signaler les échecs au lieu de les corriger en silence.
\iqlookfor{provenance, gestion des versions des retraitements, traitement en flux, rapprochement.}
\end{solution}

\begin{solution}{iq:in:reading-a-trading-firms-accounts:4}
Une bonne année pour le marché ou la stratégie~; un changement d'entité ou de prix de transfert~; des salariés transférés vers une autre société du groupe
ou reclassés comme prestataires ou associés~; un changement d'exercice comptable~; une variation de change~; un gain exceptionnel dans
le chiffre d'affaires~; une acquisition consolidée tardivement dans l'année.
\iqlookfor{effets d'entité et de période, avant la productivité.}
\end{solution}

\begin{solution}{iq:in:reading-a-trading-firms-accounts:5}
La réglementation (une seule entité détient les autorisations de trading et le capital), la forme juridique (une société de personnes se prête au partage du bénéfice
avec des associés), la fiscalité et le prix de transfert dans tout le groupe, et la possibilité de détacher du personnel auprès de plusieurs entités du groupe.
\iqlookfor{fonds propres réglementaires, économie des sociétés de personnes, prix de transfert.}
\end{solution}

\begin{solution}{iq:in:reading-a-trading-firms-accounts:6}
Régresser les \omterm{def:in:reading-a-trading-firms-accounts:staff}{charges de personnel} sur le chiffre d'affaires sur plusieurs années (livre~16, chapitre~1)~: la pente estime la part variable par unité de
chiffre d'affaires, l'ordonnée à l'origine la part fixe. Six observations donnent un intervalle large~; la croissance des effectifs, les rémunérations différées passées en charges sur
plusieurs années et les coûts ponctuels la biaisent.
\iqlookfor{la régression, son incertitude avec peu d'observations, et le calendrier de la rémunération différée.}
\end{solution}
